%%%%%%%%%%%%%%%%%%%%%%%%%%%%%%%%%
% Conclusion
%%%%%%%%%%%%%%%%%%%%%%%%%%%%%%%%%

\section{Conclusion}
\label{sec:conclusion_cm}

We have derived the probability distributions, closed-form expected values, and variances of the principal binary classification measures---positive predictive value, negative predictive value, accuracy, the $\fone$-score, and Youden's $J$ statistic---for the Bernoulli set model.
In each case the measure is a random variable whose distribution follows from the binomial error counts established by the element-wise independence axiom (\cref{asm:element_indep}).
The expected PPV and $\fone$-score were obtained via second-order delta method expansions (\cref{thm:approx_expected_precision,thm:f1_approx}), and their accuracy was validated by Monte Carlo simulation (\cref{fig:ppv_theory_vs_sim,fig:f1_theory_vs_sim}).

These results apply to any data structure satisfying the Bernoulli axioms, including Bloom filters~\cite{bf} and perfect hash filters~\cite{phf}.
In particular, the sensitivity of PPV to prevalence (\cref{fig:ppv_vs_fprate}) has practical implications for system design: a data structure with a nominally small false positive rate may still produce unreliable positive predictions when the base rate of positives is low.
Designers should therefore evaluate their systems not only by FPR and FNR but also by the predictive values that the intended query workload will induce.

Several directions remain open.
First, the distribution of the Matthews correlation coefficient (MCC) under the Bernoulli model has not been derived in closed form; it involves a more complex ratio of binomial random variables than PPV or $\fone$.
Second, extension to correlated error models---where the element-wise independence axiom is relaxed---would broaden the applicability of the framework to data structures with dependent hash functions or adversarial inputs.
Finally, the interplay between prevalence and the higher-order composition algebra of~\cite{bernoulliSets,bernoulliComposition} deserves further study: as approximate sets are composed through union and intersection, the effective prevalence at each stage changes, and so do the classification measures.
